\subsection{Higher-order pose remainder diagnostic}\label{sec:cubicpose}
The enclosing-domain calculation extends to Taylor degree $d$. Let $S_j$ be
the $L^2$ Frobenius norm of the known Fourier adjoint field's $j$th derivative
tensor over the enclosing domain. For the rigid pose path
$y(t)=\exp(t\Omega)x+t b$, use
$M_1=\operatorname{hypot}(a\sqrt{3}/2,s_t\|B_{\rm det}\|)$ and
$M_j=a^j\sqrt{3}/2$ for $j\ge2$, where $B_{\rm det}$ maps detector shifts to
object coordinates. The joint pose ball justifies the hypotenuse expression.
The chain rule, tensor Cauchy--Schwarz and volume preservation give
\[
 \left\|(f\circ y)^{(d+1)}\right\|_{L^2}
 \le \sum_{j=1}^{d+1}\mathcal B_{d+1,j}(M_1,\ldots,M_{d+1})S_j,
\]
with partial exponential Bell polynomials $\mathcal B$. Division by $(d+1)!$
bounds the Taylor remainder. The existing approximate-embedding phase residual
also needs a higher-order Leibniz pad. These are standard differentiation and
norm inequalities, not a new general statistical theorem.

A declared two-case feasibility probe retains every degree in
Table~\ref{tab:higherremainder}, at identical weights, acquisition geometry and
one-degree/0.5\,\AA\ joint pose budgets. Both use EMPIAR-10049 geometry and the
same $B=2,P=1$ density class. The values include multiplication by $B+P$.
They show a potential reduction in one bias component. They are \emph{not}
confidence-interval widths: replacing a remainder while omitting the extra
polynomial fields would be invalid. The first feature's degree-two enclosing
remainder is 3.550 of its total bias bound 5.155; its sign power remains
negligible. No degree was selected by comparing resulting coverage.
The Gaussian sigmas are respectively 20 and 10\,\AA; both use frequency radius
12.

\begin{table}[h]
\centering
\caption{Remainder-only bounds at the same weights. Degree means retained Taylor
polynomial degree, not derivative order. No higher-order interval is inferred
from this table.}\label{tab:higherremainder}
\begin{tabular}{lrrrr}
\toprule
Feature / particles & $d=2$ & $d=3$ & $d=4$ & $d=5$\\
\midrule
Pilot region, 128 & 3.550 & 0.815 & 0.160 & 0.027\\
Center, 1,024 & 21.110 & 5.694 & 1.287 & 0.250\\
\bottomrule
\end{tabular}
\end{table}

Analytic coefficient identities and independent nonlinear Taylor checks
support the implementation. Floating-point guards remain heuristic.

\paragraph{Full cubic extension.}
A separate bounded diagnostic retains all cubic fields for the first case;
its protocol and code precede its empirical outcome. Write the nonlinear phase
as $\Phi_0+t\Phi_1+t^2\Phi_2+t^3\Phi_3+\cdots$. The third exponential
coefficient is
\[
\begin{aligned}
 E_3&=i\Phi_3-\Phi_1\Phi_2-i\Phi_1^3/6,\\
 \Phi_3&=-a^2\|u_{\rm rot}\|^2\Phi_{1,\rm rot}/6.
\end{aligned}
\]
The last identity uses $[u]_\times^3=-\|u\|^2[u]_\times$; dropping it would
omit rotation curvature. For degree $m\in\{1,2,3\}$, the normalized symmetric
lift $z_\alpha=\sqrt{m!/\alpha!}\,u^\alpha$, $|\alpha|=m$, has
$\|z\|^2=\|u\|^{2m}\le1$. Thus there are $5+15+35=55$ columns per particle,
using twenty spatial monomials. Positive degree/particle scales $d_{im}$ give
a joint lifted radius $L=\sqrt{\sum_{im}d_{im}}$. Cheap quadrature chooses
these scales only; order-80 integration and a degree-six product-kernel pad
protect the continuous spectral audit. Four fresh Gaussian probes and forty
power steps use failure probability $10^{-6}/12$.

For the pilot-corrected estimator in Equation~\eqref{eq:estimator},
$\widehat\theta=\ip{\ell}{\rho_0}+w^\top(y-A_0\rho_0)$,
the residual cross-term argument remains valid with this larger field operator.
For $H\ge\|\ell-A_0^*w\|$, $c\ge\|F^*(\ell-A_0^*w)\|$,
$f\ge\sup_{\|v\|\le L}\|Fv\|$, $p\ge\|F^*\rho_0\|$ and quartic field
remainder $r$, the conditional bias bound is
\[
 B\sqrt{H^2+2Lc+f^2}+\min(Pf,Lp)+(B+P)r.
\]
Its proof expands the squared residual norm and applies Cauchy--Schwarz to the
known pilot. Twenty continuous cell and Gaussian moments evaluate the pairings.
Independent scalar phase recurrences, direct quadrature, adjoint identities and
the old quadratic restriction check signs and normalization. The quadratic and cubic intervals are alternatives; their unadjusted minimum
is not a joint confidence claim. The empirical
cubic outcome is pending at this development checkpoint. A correct extension
alone does not establish practical usefulness or calibrated pose budgets.
